\originalchapter{蛋白折叠与能量景观}
\originalsection{构象数和随机能量模型}
氨基酸序列确定相互作用，但一条链有指数多种构象。低能构象是否能在热涨落中保持优势，取决于能量与构象熵的竞争。随机能量模型把不同紧密构象的能量视为独立高斯变量。这一独立性忽略真实构象间相关性，却能分清“有最低能量”和“能稳定折叠到特定结构”两件事。

设构象数 $\Omega=\e^{Ns_0/\kb}$，能量方差 $\Sigma^2=Nu^2$。能量 $E$ 附近的平均构象数正比于 $\Omega\e^{-E^2/(2\Sigma^2)}$，所以微正则熵为 $S(E)=\kb\ln\Omega-\kb E^2/(2\Sigma^2)$，只在 $S\ge0$ 的区间可信。低能截断为 $E_c=-\Sigma\sqrt{2\ln\Omega}$。以 $F=E-TS$ 最小化，鞍点 $E_*(T)=-\Sigma^2/(\kb T)$；当它越过截断，真实能量停在 $E_c$，不能继续把负的构象熵当作实际态数。

冻结温度为 $\kb T_c=\Sigma/\sqrt{2\ln\Omega}$。高温自由能 $F_{\rm mis}=-\kb T\ln\Omega-\Sigma^2/(2\kb T)$，低温为 $E_c$。对一个设计良好的序列，另有唯一的天然态 $E_N<E_c$。比较 $E_N$ 与 $F_{\rm mis}$ 能确定折叠温度：过低的天然能量开出能隙，降低误折叠态竞争。序列设计使用的“温度”是挑选序列的统计参数，不是细胞的实际温度。

\originalsection{折叠动力学和势垒}
自由能最低并不保证快速抵达。以构象坐标 $q$ 的有效自由能 $G(q)$ 描述折叠，过阻尼随机运动满足 $\dot q=-M G'(q)+\sqrt{2M\kb T}\,\eta(t)$，其中 $M$ 为迁移率，白噪声满足 $\avg{\eta(t)\eta(t')}=\delta(t-t')$。相应概率流是 $J=-M(G'p+\kb T p')$；零流分布为 $p\propto\e^{-\beta G}$，这也核对噪声强度。

对高势垒 $\Delta G\gg\kb T$，逃逸时间的主指数为 $\e^{\beta\Delta G}$。指数可从准平衡势阱权重与势垒附近稀有概率的比值看出；更精确的 Kramers 前因子还依赖势阱和势垒曲率及阻尼机制。不能仅据一个自由能差写出绝对折叠时间。多个中间态、陷阱和并行路径会破坏单指数弛豫。

一个可检验的两态模型令未折叠和已折叠概率分别为 $p_U,p_F$，速率 $k_f,k_u$。$\dot p_F=k_f(1-p_F)-k_up_F$，弛豫率为 $k_f+k_u$，平衡折叠比例 $k_f/(k_f+k_u)$；详细平衡要求 $k_f/k_u=\e^{-\beta(G_F-G_U)}$。平衡比例能约束速率比，不能独立确定两个速率。这里的动力学解释是对课程主题的编者补充，不冒称缺少详细PDF的折叠讲次逐页翻译。
